\documentclass{ximera}[font=12pt]


\title{Absolute Value and Opposites}   % Use xmlatex all -s to compile when SERVE refuses to work.
\author{Kelly Stady}
\license{CC: 0}         % replace with an appropriate license, or set it in xmPreamble

\begin{document}   

\begin{abstract}

\end{abstract}

\maketitle

\label{xim:AbsValue}


\section*{Absolute Value}

If you're cutting a piece of wood to be 10 inches in length, being 0.5 inches too long is just as problematic as being 0.5 inches too short. In both cases, the \textbf{size} 
of the error is 0.5 and this is what is most important.  The concept of \textbf{absolute value} can be used to represent the \textbf{magnitude} of a mistake, regardless 
of whether it was an overshot or undershot.

The \textbf{absolute value of a number} is its \textbf{distance from zero} on a number line.  Since distance is always positive or 0, so is absolute value.  Absolute value is represented by two vertical bars and can be interpreted as the \textbf{size}, \textbf{strength} or \textbf{magnitude} of a number.

The number line below shows the absolute values of $-4$ and $4.$  Since their absolute values are the same, the numbers $-4$ and $4$ have the same strength or magnitude.

\begin{center}
\begin{tikzpicture}[scale=1.2, font=\sffamily]
% Draw Number Line
\draw[<->, thick] (-5.5,0) -- (5.5,0);
\foreach \x in {-5,-4,...,5}
    \draw (\x, 0.1) -- (\x, -0.1) node[below] {\small \x};

% Highlighting the points -4, 0, and 4
\filldraw[orange!85!black] (-4,0) circle (2pt);
\filldraw[glaucous] (4,0) circle (2pt);
\filldraw[black] (0,0) circle (2pt);

% Horizontal line for -4 (Red) - Original position
\draw[orange!85!black, ultra thick, |-|] (-4, 0.6) -- (0, 0.6)
    node[midway, above, font=\bfseries] {Distance = 4};

% Horizontal line for +4 (Blue) - Original position
\draw[glaucous, ultra thick, |-|] (0, 1.3) -- (4, 1.3)
    node[midway, above, font=\bfseries] {Distance = 4};

% NEW: Absolute values centered UNDER the number line
\node[orange!85!black, font=\bfseries] at (-2, -0.8) {$\mathbf{|-4| = 4}$};
\node[glaucous, font=\bfseries] at (2, -0.8) {$\mathbf{|4| = 4}$};

% Vertical guide lines - Original position
%\draw[dashed, gray] (0,0) -- (0, 1.5);
%\draw[dashed, gray] (-4,0) -- (-4, 0.6);
%\draw[dashed, gray] (4,0) -- (4, 1.3);
\end{tikzpicture}
\end{center}

\textbf{To find the absolute value of a number}, think of the number line as a path and each integer as a step. The 
absolute value of a number is how many steps you would have to take to travel from the number to zero.  For example, 
you would have to take 2 steps to travel from $-2$ to $0,$ so the absolute value of $-2$ is $2.$
\[ |-2| = 2 \]


\begin{problem}
Find the absolute value of each number.

\[ |-1| = \begin{prompt} \answer{1} \end{prompt} \]

\[ |7| = \begin{prompt} \answer{7} \end{prompt} \]


\[ |-55| = \begin{prompt} \answer{55} \end{prompt} \]

\end{problem}



\begin{problem}
You are building a custom bookshelf and need to cut boards that are $36$ inches long.  A board will be usable if its 
length is off by no more than $0.25$ inch of the target.  Suppose you cut the first board $35.8$ inches long.  The 
difference between the actual board length and the target length is the error.
\begin{eqnarray*}
  \textbf{Error} &=& 35.8 - 36 \\
        & & \\
        &=& -0.2 \  \textrm{inches}
\end{eqnarray*}

The negative sign indicates that you cut the board too short.  However, it only matters if the board is useable or not, so we 
take the \textbf{absolute value} to find the \textbf{distance} from the target length.  This distance is called the \textbf{absolute error}.

~ \\[0.5em]


\begin{eqnarray*}
\textbf{Absolute Error} &=& |35.8 - 36| \\
    & & \\
    &=& |-0.2| \\
    & & \\
    &=&  \answer{0.2} \ \textrm{inches}
\end{eqnarray*}

Is this board usable? \textbf{Hint}: A board is usable if the absolute error is less than or equal to $0.25$ inches.

\begin{multipleChoice}
    \choice[correct]{Yes, the absolute error is less than $0.25$ inches.}
    \choice{No, the absolute error is greater than $0.25$ inches.}
\end{multipleChoice}


\end{problem}


\newpage

~ \\[2em]

\insightIcon

\section*{Extra Insight}

In physics, the motion of an object is often described by its \textbf{velocity}, which indicates its speed \emph{and} direction.  
When an object moves \textbf{upward} or to the \textbf{right}, its velocity is \textbf{positive} and when it moves 
\textbf{downward} or to the \textbf{left}, its velocity is \textbf{negative}.  We use
absolute value to turn velocity into speed.

\begin{center}
Speed = |Velocity|  \\

\textbf{\textcolor{glaucous}{Speed is the magnitude or size of velocity}}. 
\end{center}


When a ball is thrown straight up in the air, its velocity is positive as it moves upward, zero when it reaches its highest point and negative when
it falls back to Earth.  If, after 6 seconds, the velocity of the ball is $-15$ feet per second, we know that it is falling back to Earth,
since its velocity is negative.  Since $|-15| = 15$ the ball is traveling at a speed of 15 feet per second 6 seconds after it was thrown upward.


\begin{center}
\begin{tikzpicture}[>=stealth, scale=1.2]

    % Ground line
    \draw[ultra thick] (-1.5, 0) -- (4.5, 0);
    \fill[pattern=north east lines] (-1.5, -0.2) rectangle (4.5, 0);

    % The Ball (Added at the start of the blue line)
    %\shade[ball color=magenta] (0.8, 0.5) circle (0.15);  Alternative for the ball
    \shade[ball color=white, draw=black, thin] (0.8, 0.5) circle (0.15);

    % Upward Path (Positive Velocity)
    \draw[->, glaucous!70!black, line width=3pt] (0.8, 0.7) -- (0.8, 3.5);
    \node[left, glaucous!70!black, font=\bfseries] at (0.8, 2) {\textbf{Velocity (+)}};

    % The Peak (Zero Velocity)
    \draw[dashed, thick] (-1, 3.5) -- (4, 3.5);
    \node[right] at (4, 3.5) {\textbf{Velocity = 0}};

    % Downward Path (Negative Velocity)
    \draw[->, orange!80!black, line width=3pt] (2.2, 3.5) -- (2.2, 0);
    \node[right, orange!80!black, font=\bfseries] at (2.2, 2) {Velocity  (--)};


\end{tikzpicture}
\end{center}



\begin{problem}
A model rocket was launched straight up in the air.  After 8 seconds, the velocity of the rocket was recorded as $-12$ feet per second.
In which direction was the rocket moving 8 seconds after it was launched?

~ \\[2em]

\begin{multipleChoice}
\choice[correct]{Downward}
\choice{Upward}
\end{multipleChoice}

~ \\[2em]

What was the rocket's speed at that moment?

Speed = $\answer{12}$ feet per second

\end{problem}

~ \\[1em]

\section*{Opposites}

~ \\[1em]

\begin{definition}[\textbf{Opposites}]

~ \\[1em]

Two numbers that are the same distance from zero but on opposite sides
of zero are called \textbf{opposites}.

The opposite of a number is found by changing its sign.

The phrase ``\textbf{the opposite of}" is written mathematically as a dash ($-$).

\end{definition}



\begin{example}
The numbers 3 and $-3$ are opposites.  Notice that they lie the same distance from zero on the number line, but in \textbf{opposite} directions, and they have the same absolute values.

~ \\[2em]

\begin{center}
\begin{tikzpicture}[scale=1.2, font=\sffamily]
% Draw Number Line
\draw[<->, thick] (-4.5,0) -- (4.5,0);
\foreach \x in {-4,...,4}
    \draw (\x, 0.1) -- (\x, -0.1) node[below] {\small \x};

% Highlighting the points -3, 0, and 3
\filldraw[orange!85!black] (-3,0) circle (2pt);
\filldraw[glaucous] (3,0) circle (2pt);
\filldraw[black] (0,0) circle (2pt);

% Horizontal line for -3 (Orange)
\draw[orange!85!black, ultra thick, |-|] (-3, 0.6) -- (0, 0.6)  % |-| is a line with two vertical bars at its end
    node[midway, above, font=\bfseries] {Distance = 3};

% Horizontal line for 3 (Blue)
\draw[glaucous, ultra thick, |-|] (0, 1.3) -- (3, 1.3)
    node[midway, above, font=\bfseries] {Distance = 3};

% Absolute values- centered
\node[orange!85!black, font=\bfseries] at (-1.5, -0.8) {$\mathbf{|-3| = 3}$};
\node[glaucous, font=\bfseries] at (1.5, -0.8) {$\mathbf{|3| = 3}$};

\end{tikzpicture}
\end{center}
\end{example}

~ \\[1em]

\subsection*{Properties of Opposites}

\begin{enumerate}
\item Opposites have the same absolute value because they are equidistant from zero.

\item One is positive and the other is negative.

\item Adding a number and its opposite always results in zero: $a +(-a) = 0.$

\item Zero is unique because it is its own opposite: $0= -0.$

\end{enumerate}

~ \\[1em]

\begin{problem}
Complete the table.

~ \\[1em]

\[
\renewcommand{\arraystretch}{1.8}
\begin{array}{c @{\hspace{30pt}} c @{\hspace{30pt}} c}
\text{Number} & \text{Opposite} & \text{Absolute Value} \\
\hline
-15 & \answer{15} & \answer{15} \\
8 & \answer{-8} & \answer{8} \\
0 & \answer{0} & \answer{0} \\
-2.5 & \answer{2.5} & \answer{2.5} \\
\end{array}
\]

\end{problem}


\subsection*{The Three Meanings of a Dash ($-$)}

In mathematics, the dash is used in three different ways depending on its role.

\begin{enumerate}
    \item \textbf{Negative} \ When a dash is attached directly to a number, it describes its location on the number line and
    is \textbf{read} as \textbf{negative}. \\

    \item \textbf{Opposite} \ A dash placed before parentheses, absolute value or a variable represents an action: changing the sign or multiplying by 
    $-1$ and is \textbf{read} as \textbf{the opposite of}. \\

    \item \textbf{Subtraction} \ A dash placed between two numbers represents the operation of subtraction and is \textbf{read} as \textbf{minus}.
\end{enumerate}


~ \\[1em] 

\begin{example}

\renewcommand{\arraystretch}{1.5}
\begin{tabular}{p{0.15\textwidth} p{0.5\textwidth}}
\large \textbf{\textcolor{glaucous}{Expression}} & \large \textbf{\textcolor{glaucous}{How to Read It}} \\
 $-8$ & ``Negative 8'' \\
$-(-8)$ & ``The opposite of negative 8'' \\
& $-(5)$ & ``The opposite of 5'' \\
& $-|6|$ & ``The opposite of the absolute value of $6$'' \\
& $-x$ & ``The opposite of $x$'' \\
& $10 - 8$ & ``10 minus 8'' \\
\end{tabular}

\end{example}


\begin{remark}
The expression $-x$ is sometimes read as negative $x.$  However, keep in mind that the expression
could be positive, negative or even 0, depending on the value of $x.$  
\end{remark}


\begin{problem}
Identify the correct way to read each expression.

~ \\[2em]

\begin{center}
\renewcommand{\arraystretch}{1.5}
\begin{tabular}{cl}
\textbf{Expression} & \textbf{How to Read It} \\ 

$18 - 5$ & \wordChoice{\choice{Negative 5} \choice[correct]{Eighteen minus five} \choice{The opposite of 5}} \\ 

$-(-12)$ & \wordChoice{\choice{Negative negative 12} \choice{Minus negative 12} \choice[correct]{The opposite of negative 12}} \\ 

$-20$ & \wordChoice{\choice[correct]{Negative 20} \choice{The opposite of negative 20} \choice{Minus 20}} \\ 


$-|-7|$ & \wordChoice{\choice[correct]{The opposite of the absolute value of negative 7} \choice{The opposite of the absolute value of 7} \choice{The absolute value
of negative 7}} \\ 

$-k$ & \wordChoice{\choice{Dash k} \choice{Minus k} \choice[correct]{The opposite of k}} \\ 

\end{tabular}
\end{center}
\end{problem}



\begin{definition}[\textbf{The Double Negative Property}]

For any number $a,$

\[ -(-a) = a \]

The expression $-(-a)$ is read as ``the opposite of negative $a$."


\begin{problem}
Use the Double Negative Property to simplify the following expressions.

~ \\[2em]

\begin{enumerate}
\item $-(-14) = \answer{14}$
\item $-(-100) = \answer{100}$
\item $-(-0) = \answer{0}$
\end{enumerate}
\end{problem}

\end{definition}


\begin{problem}
How is the expression $-(-22)$ read?

~ \\[2em]
\begin{multipleChoice}
\choice{Negative negative 22}
\choice{Minus negative 22}
\choice[correct]{The opposite of negative 22}
\choice{The opposite of 22}
\end{multipleChoice}

\end{problem}


\begin{problem}
If $x = -10$, what is the value of $-(-x)$?

~ \\[2em]

\begin{multipleChoice}
    \choice{$10$}
    \choice[correct]{$-10$}
    \choice{$0$}
\end{multipleChoice}

\begin{feedback}
Remember: $-(-x) = x,$ so since $x$ started as $-10$, the final result is also $-10$.
\end{feedback}
\end{problem}


\begin{warning}
A common mistake is to treat absolute value like parentheses. Keep in mind that \textbf{parentheses} are used for \textbf{grouping}, whereas \textbf{absolute value} is an \textbf{operation}.


\begin{description}[style=nextline, leftmargin=0pt]  %leftmargin removes the indentation below items
    \item[Case 1: Double Negative with Parentheses]

Since we have parentheses, we simply find the opposite of negative 7, which is 7.
\[ -(-7) = 7 \]

\item[Case 2: Double Negative with Absolute Value]

To find the opposite of the absolute value of negative 7, start with the absolute value: $|-7| = 7,$ then find the opposite of that result.  The opposite of 7 is $-7$
therefore,

\[ -|-7| = -7. \]

\end{description}

\end{warning}

\subsection*{Absolute Value vs. Opposite}

While absolute value always results in a nonnegative (zero or positive) number, taking the opposite of a number flips its sign. Therefore, an opposite can be positive, negative or 0. Note that for the number 0, both its absolute value and its opposite are 0.

Below, the concept of an opposite has been used to write a rule for finding absolute value.  Notice that this rule doesn't rely on distance.

\begin{procedure}[\textbf{Finding the Absolute Value of a Number}]
\begin{description}[style=nextline, leftmargin=0pt, font=\bfseries]
    \item[Positive Numbers and 0]
    The absolute value of a positive number or 0 is the number itself.
    \item[Negative Numbers]
    The absolute value of a negative number is the \textbf{opposite} of that number.
\end{description}
\end{procedure}

Later, you will see how these two concepts are used together to add integers.

\begin{problem}
\textbf{True or False}: If a number is negative, its opposite must be positive.

~ \\[2em]

\begin{multipleChoice}
\choice[correct]{True}
\choice{False}
\end{multipleChoice}


\begin{feedback}[correct]
Exactly! The opposite flips the sign, so a negative becomes positive.
\end{feedback}

\end{problem}

\begin{problem}
\textbf{True or False}: Two different numbers can have the same absolute value.

~ \\[2em]

\begin{multipleChoice}
\choice[correct]{True}
\choice{False}
\end{multipleChoice}

\begin{feedback}[correct]
Yep! The numbers $2$ and $-2$ have the same absolute values.  In fact, all opposites have the same absolute values.
\end{feedback}

\end{problem}

\begin{problem}
Which has a greater value: the opposite of $-10$ or the absolute value of $-10?$  Explain.

~ \\[2em]

\begin{multipleChoice}
\choice[correct]{Neither, they are the same.}
\choice{The opposite of $-10.$}
\choice{The absolute value of $-10.$}
\end{multipleChoice}

~ \\[2em]

\begin{feedback}
The opposite of $-10$ is 10 and the absolute value of $-10$ is also 10, so they are the same.
\end{feedback}

\end{problem}


\begin{problem}
Two numbers, $a$ and $b$ are opposites. If the distance between
$a$ and $b$ on a number line is 10 units, what are the two numbers?

\begin{multipleChoice}
\choice{$-10$ and $10$}
\choice[correct]{$-5$ and $5$}
\choice{$-6$ and $4$}
\choice{$-7$ and $3$}
\end{multipleChoice}

\begin{hint}
Draw a number line and with $a$ and $b$ ten units apart.
\end{hint}

~ \\[2em]

\begin{feedback}
Good job!  You made it through this problem and this section.
\end{feedback}

\end{problem}



\end{document}

